\documentclass[11pt]{article}
\usepackage{coursenotes}
\begin{document}
\handoutheader{H2}{Neural and foundation-model estimators}{What deep learning adds to effect estimation, and what it cannot}{Meetings 3 (Conditional effects) and 9 (Generative AI)}

\begin{decision}
A vendor offers a deep-learning model that ``estimates each customer's response from their app behavior and
reviews, and corrects for unobserved confounders''. Before buying it, a firm needs to know which parts of that
claim are about estimation (where neural networks can help) and which are about identification (where no
architecture can help without assumptions).
\end{decision}

\begin{goals}
  \item describe TARNet, counterfactual regression (CFR), and DragonNet, and what each one's loss targets;
  \item explain how balancing a representation can hide an overlap failure;
  \item state why CEVAE-type models cannot recover hidden confounding without stated proxy structure;
  \item describe in-context (foundation-model) estimators and locate their identifying assumptions in the prior;
  \item evaluate a neural estimator without being misled by benchmark artifacts.
\end{goals}

%=====================================================================================
\sectionone

\subsection{Where neural networks help}

With treatment $D$, outcome $Y$, and pre-treatment covariates $X$, write $\mu_d(x) = \E[Y \mid D=d, X=x]$ for the
outcome regression in arm $d$ and $\tau(x) = \E[Y(1) - Y(0) \mid X=x]$ for the conditional average treatment effect
(CATE). Under conditional independence, $\{Y(0), Y(1)\} \indep D \mid X$, and overlap, $0 < \Pr(D=1 \mid X=x) < 1$,
the CATE is $\tau(x) = \mu_1(x) - \mu_0(x)$. The problem is then to estimate two regression functions well. Neural networks are
useful when $x$ is high-dimensional or unstructured (text, images, event sequences) and when the two arms
share structure, so that learning features from all the data helps both. They do not change what is identified.

\subsection{TARNet: a shared representation with two heads}

TARNet maps $x$ through shared layers to a representation $\Phi(x)$ and then through two heads $h_0, h_1$, giving
$\hat\mu_d(x) = h_d(\Phi(x))$. It minimizes the \emph{factual} loss
$\sum_i \big(Y_i - h_{D_i}(\Phi(X_i))\big)^2$: each unit trains only its own arm's head.

\begin{prop}{What the factual loss can learn}{factual}
With unrestricted $\Phi$ and heads, the factual loss is minimized by any functions with $h_d(\Phi(x)) = \mu_d(x)$ at
every $x$ where $\Pr(D = d \mid X = x) > 0$. At $x$ where arm $d$ is never observed, the loss does not depend on
$h_d(\Phi(x))$.
\end{prop}
\begin{proof}
The expected loss is $\sum_d \E\big[\ind{D=d}(Y - h_d(\Phi(X)))^2\big]$; for each $d$ and $x$ with positive
probability of arm $d$, squared loss is minimized by the conditional mean $\mu_d(x)$. Where $\Pr(D=d \mid x) = 0$, the
term has weight zero.
\end{proof}

Outside overlap the network's prediction is decided by architecture and regularization, not by data: it is an
extrapolation that looks like an estimate.

\subsection{Counterfactual regression: balancing the representation}

CFR adds a penalty $\alpha \cdot \text{IPM}(\Phi \mid D=1,\ \Phi \mid D=0)$ that makes the distribution of $\Phi(X)$ similar
in the two arms (an integral probability metric such as the Wasserstein distance or MMD). Shalit et al.\ (2017) bound the estimation error of $\tau$ by the factual losses plus a constant times this distance, which
motivates the penalty. The bound assumes $\Phi$ is invertible. A penalized network need not be, and then balancing has
a cost.

\begin{prop}{Balancing can hide an overlap failure}{hide}
If $\Phi$ maps a region where only treated units are observed onto the same representation as a region where
controls are observed, then $\hat\mu_0$ in the first region equals the control outcome mean of the second, whatever
the true $\mu_0$ there.
\end{prop}
\begin{proof}
$\hat\mu_0(x) = h_0(\Phi(x))$ depends on $x$ only through $\Phi(x)$; if $\Phi(x) = \Phi(x')$ then $\hat\mu_0(x) = \hat\mu_0(x')$, and
the factual loss fits $h_0$ at that representation to controls at $x'$.
\end{proof}

The penalty is smallest exactly when the network forgets the features that separate treated from control units,
which are often the confounders. Balance in $\Phi$ is then a symptom of lost information, not evidence of
comparability. Check overlap on $X$ (propensity scores), not on the learned representation.

\subsection{DragonNet: adding the propensity}

DragonNet adds a third head that predicts $D$ from $\Phi(x)$. By the balancing property of the propensity score
$e(x) = \Pr(D=1 \mid X=x)$, conditioning on $e(X)$ suffices for adjustment, so a representation need only keep the information that
predicts treatment and outcome; the propensity head pushes $\Phi$ to keep it. A \emph{targeted regularization} term adds
a one-parameter correction that makes the final estimate behave like the doubly robust (AIPW) estimator. The
architecture's gains are in estimation efficiency; identification still rests on conditional independence and
overlap.

\subsection{Hidden confounders and the causal effect VAE}

The causal effect variational autoencoder (CEVAE) posits a latent $Z$ that causes $X$, $D$, and $Y$, fits a VAE to the
observed $(X, D, Y)$, and computes effects as if $Z$ were the only confounder. The claim is that the proxies $X$ let
the model recover $Z$.

\begin{prop}{The observed distribution does not pin down the effect}{nonid}
Two models can generate the same joint distribution of $(D, Y)$ with different effects. Example: model A has no
confounder, $\Pr(D=1) = 0.5$, and $\Pr(Y=1 \mid D) = 0.4 + 0.2D$, so $\tau = 0.2$. Model B has a binary $U$ with
$\Pr(U=1) = 0.5$, $\Pr(D=1 \mid U) = 0.2 + 0.6U$, and $\Pr(Y=1 \mid U) = \tfrac13 + \tfrac13 U$ independent of $D$, so
$\tau = 0$. Both give $\Pr(Y=1 \mid D=1) = 0.6$ and $\Pr(Y=1 \mid D=0) = 0.4$.
\end{prop}
\begin{proof}
In B, $\Pr(U=1 \mid D=1) = 0.8$ and $\Pr(U=1 \mid D=0) = 0.2$, so $\Pr(Y=1 \mid D=1) = 0.8 \cdot \tfrac23 + 0.2 \cdot \tfrac13 = 0.6$
and $\Pr(Y=1 \mid D=0) = 0.2 \cdot \tfrac23 + 0.8 \cdot \tfrac13 = 0.4$.
\end{proof}

Any method fitted to the observed distribution, however flexible, cannot choose between A and B. Proxies help only
under explicit structure. \emph{Proximal causal inference} (Miao et al.\ 2018) shows
identification when two proxies play distinct roles: a negative-control \emph{exposure} (related to $U$, with no
effect on $Y$) and a negative-control \emph{outcome} (related to $U$, unaffected by $D$), plus a completeness
condition. CEVAE leaves that structure implicit, and Rissanen and Marttinen (2021) show it can be inconsistent when
proxies are weak or the latent model is misspecified. Treat a claim to ``correct for unobserved confounders'' as
an identification claim and ask for the assumptions.

\subsection{Time-varying treatments}

With treatments over time (a sequence of offers), past treatments affect future covariates that affect future
treatments: \emph{time-varying confounding}. Sequence models (recurrent networks, the causal transformer of
Melnychuk et al.\ 2022) estimate counterfactual trajectories, using balancing representations at
each step. The identifying assumption is sequential conditional independence (at each step, treatment is independent of
the potential outcomes given the history of past treatments and covariates). It is the same assumption
G-computation needs, and G-computation remains the transparent baseline.

\subsection{In-context and foundation-model estimators}

A \emph{prior-fitted network} (PFN) is a transformer trained on millions of simulated datasets, each drawn from a prior
over structural causal models. Given a new dataset in its context window, it outputs an estimate of an interventional
quantity directly, approximating the posterior-predictive answer under that prior. Do-PFN (Robertson et al.\ 2025)
applies this to effect estimation.

\begin{prop}{The identification lives in the prior}{pfn}
If two causal models in the prior's support generate the same observed distribution but different effects, the
data cannot separate them, and an ideal PFN returns a prior-weighted combination of their effects.
\end{prop}
\begin{proof}
The posterior over models is proportional to prior times likelihood; equal likelihoods leave the prior's relative
weights unchanged, so the posterior-predictive effect averages the two effects with prior weights.
\end{proof}

When the effect is identified (for example, conditional independence holds in every model of the prior), a PFN can
be an efficient estimator with no tuning. When it is not, its answer reflects the simulator designer's prior over
worlds, which is where its assumptions must be examined.

\subsection{Evaluating neural estimators}

Semi-synthetic benchmarks (IHDP, ACIC) simulate outcomes on real covariates so the true effect is known. Curth et al.\
(2021) show that rankings on them can reflect how the outcomes were simulated rather than real-world performance; for
example, benchmarks whose effects are simple functions of features that the response surfaces also use reward
architectures that share structure across arms. Evaluate on the decision instead: the value of the targeting policy a
model implies, on held-out randomized data, and a test of whether its estimated effects are calibrated on
the same data.

\begin{pitfall}[reading a deep-learning causal claim]
Treating representation balance as overlap; treating a VAE's latent variable as a measured confounder; ranking models
on one semi-synthetic benchmark; accepting a foundation model's estimate without asking what its prior assumes.
\end{pitfall}

%=====================================================================================
\sectiontwo

\paragraph{Setting.} QuickBite wants to estimate the effect of a ``featured restaurant'' slot on a restaurant's weekly
orders from two years of logs. Restaurants are described by menu text, photos, and history. Three regions of covariate
space matter: (a) small new restaurants, (b) mid-sized ones, (c) large chains. The slot was sold mostly to chains, and in
region (c) every restaurant was featured.

\paragraph{Step 1: overlap.} A propensity model on menu embeddings and history gives $\hat e \approx 0.15$ in (a), $0.5$ in
(b), and $1.00$ in (c). No unfeatured chain exists, so $\tau$ is not identified in (c) (Result~\ref{prop:factual}).

\paragraph{Step 2: what CFR does.} A CFR model with a strong balancing penalty reports small imbalance and an effect for
chains of $0.40$ (in thousands of orders). Inspecting $\Phi$ shows that chains and small new restaurants share a region of
representation space. True control outcomes are $0.2$ for (a) and $0.5$ for (c), and the true chain effect is $0.1$:
\[
  \hatt(c) = \hat\mu_1(c) - \hat\mu_0(c) = 0.6 - 0.2 = 0.4, \qquad \tau(c) = 0.6 - 0.5 = 0.1 .
\]
The balance came from mapping chains onto small restaurants (Result~\ref{prop:hide}), which imported the wrong
baseline.

\paragraph{Step 3: what the data support.} Report effects for (a) and (b), where overlap holds, and state that the logs say
nothing about chains. To learn the chain effect, randomize the slot among chains for a few weeks.

\begin{exercise}[computation]
In Result~\ref{prop:nonid}, change model B so that $\Pr(D=1 \mid U) = 0.3 + 0.4U$. Find $\Pr(Y=1 \mid U)$ (as $a_0 + (a_1 - a_0)U$)
that reproduces $\Pr(Y=1 \mid D=1) = 0.6$ and $\Pr(Y=1 \mid D=0) = 0.4$ with no effect of $D$. Is it a valid probability?
\end{exercise}
\answer{$\Pr(U=1 \mid D=1) = 0.35/0.5 = 0.7$ and $\Pr(U=1 \mid D=0) = 0.15/0.5 = 0.3$. Solve $0.7a_1 + 0.3a_0 = 0.6$ and
$0.3a_1 + 0.7a_0 = 0.4$: $a_1 - a_0 = 0.5$ and $a_0 = 0.25$, so $a_1 = 0.75$. Both lie in $[0,1]$, so a confounder of this
strength explains the whole association. A weaker link between $U$ and $D$ needs a stronger $U$--$Y$ link.}

\begin{exercise}[judgment]
The vendor says its foundation model ``was validated on 200 benchmark datasets with known effects''. List three
questions you would ask before using its estimates for QuickBite's slot pricing.
\end{exercise}
\answer{(1) Were the benchmark datasets generated by a prior similar to the vendor's training prior (if so, validation is
circular)? (2) Does the model's prior assume no unmeasured confounding, and is that plausible for slots sold by a sales
team? (3) How does it behave where overlap fails, such as chains: does it flag non-identification or extrapolate? Also:
can it be validated against a QuickBite randomized test?}

%=====================================================================================
\sectionthree

\subsection*{Annotated reading}
\begin{readinglist}
\reading{Shalit, Uri, Fredrik D. Johansson, and David Sontag. 2017. ``Estimating Individual Treatment Effect:
Generalization Bounds and Algorithms.'' In \emph{Proceedings of the 34th International Conference on Machine
Learning}, PMLR 70: 3076--85.}{TARNet and CFR, with the error bound that motivates balancing}{Sections 1--4}{3}
\reading{Shi, Claudia, David M. Blei, and Victor Veitch. 2019. ``Adapting Neural Networks for the Estimation of
Treatment Effects.'' In \emph{Advances in Neural Information Processing Systems 32} (NeurIPS 2019).}{DragonNet and
targeted regularization}{Sections 1--3}{3}
\reading{Louizos, Christos, Uri Shalit, Joris M. Mooij, David Sontag, Richard Zemel, and Max Welling. 2017.
``Causal Effect Inference with Deep Latent-Variable Models.'' In \emph{Advances in Neural Information Processing
Systems 30} (NeurIPS 2017).}{The CEVAE proposal}{Sections 1--3}{3}
\reading{Rissanen, Severi, and Pekka Marttinen. 2021. ``A Critical Look at the Consistency of Causal Estimation
with Deep Latent Variable Models.'' In \emph{Advances in Neural Information Processing Systems 34} (NeurIPS
2021).}{When CEVAE fails, with simulations}{Sections 1--4}{2}

\reading{Miao, Wang, Zhi Geng, and Eric J. Tchetgen Tchetgen. 2018. ``Identifying Causal Effects with Proxy
Variables of an Unmeasured Confounder.'' \emph{Biometrika} 105 (4): 987--93.
\url{https://doi.org/10.1093/biomet/asy038}}{The conditions under which proxies do identify effects}{Sections
1--3}{3}
\reading{Curth, Alicia, David Svensson, James Weatherall, and Mihaela van der Schaar. 2021. ``Really Doing Great
at Estimating CATE? A Critical Look at ML Benchmarking Practices in Treatment Effect Estimation.'' In
\emph{Proceedings of the Neural Information Processing Systems Track on Datasets and Benchmarks} (NeurIPS
2021).}{Why benchmark rankings mislead}{All}{2}
\reading{Melnychuk, Valentyn, Dennis Frauen, and Stefan Feuerriegel. 2022. ``Causal Transformer for Estimating
Counterfactual Outcomes.'' In \emph{Proceedings of the 39th International Conference on Machine Learning}, PMLR
162: 15293--329.}{Sequence models for time-varying treatments}{Sections 1--4}{3}
\reading{Robertson, Jake, Arik Reuter, Siyuan Guo, Noah Hollmann, Frank Hutter, and Bernhard Sch\"olkopf. 2025.
``Do-PFN: In-Context Learning for Causal Effect Estimation.'' In \emph{Advances in Neural Information Processing
Systems 38} (NeurIPS 2025). arXiv:2506.06039.}{A foundation-model estimator trained on simulated causal
worlds}{Sections 1--4}{3}
\end{readinglist}

\end{document}
